%-------------------------------------------------------------------------------
%	ABSTRACT PAGE
%-------------------------------------------------------------------------------

\addchaptertocentry{Хураангуй} % Хураангуйг гарчигт нэмэх

\begin{center}
{\scshape\Large \univname\par} % Их сургуулийн нэр
{\scshape\large \facname\par}\vspace{0.5cm} % Их сургуулийн нэр
{\huge\textbf{{Хураангуй}} \par}
\bigskip
{\Large{\ttitle} \par} % Тезисийн нэр
\bigskip

{\normalsize \shortname \par} % Зохиогчийн нэр
\addressname
\end{center}

\textit{\textbf{Түлхүүр үгс: \keywordnames}}
\bigskip

Хүйсийн болон нийгмийн тэгш байдалтай холбоотой асуудлууд нийгмийн олон түвшинд ажиглагдсаар байна. Тухайлбал өсвөр насны хүүхдүүдийн дунд нэгнээ ялгаварлан гадуурхах, дээрэлхэх, хөгжлийн бэрхшээлтэй иргэд амьдрах, ажиллах орчин дутмаг байхын зэрэгцээ ажлын байр, боловсролын орчин болон нийгмийн бусад хүрээнд тэгш бус харилцаа байсаар байна. Ийм асуудлууд нь зөвхөн хувь хүний сэтгэлзүй, харилцаанд нөлөөлөөд зогсохгүй нийгмийн оролцоо, тэгш боломж, хүртээмжтэй орчны хөгжилд сөргөөр нөлөөлдөг.

Эдгээр нөхцөлд иргэд эрх тэгш байдал, харилцан хүндлэл, хүртээмжтэй орчны талаарх зөв ойлголттой байх нь чухал ач холбогдолтой. Гэвч энэ төрлийн мэдээлэл нь ихэвчлэн олон эх сурвалжид тархсан, урт хэлбэртэй, заримдаа мэргэжлийн хэллэгтэй баримт бичигт агуулагддаг тул хэрэглэгчид шаардлагатай мэдээллээ богино хугацаанд олж авах, зөв ойлгох, хэрэглэхэд хүндрэлтэй байдаг.

Нөгөө талаас сүүлийн жилүүдэд байгалийн хэл боловсруулалт(NLP) болон том хэлний загвар(LLM)-т суурилсан хиймэл оюун ухааны шийдлүүд өргөн хөгжиж олон салбарт амжилттай ашиглагдаж байна. Ийм боломжийг хүйсийн болон нийгмийн тэгш, хүртээмжтэй байдлын талаарх мэдээлэл түгээхэд ашигласнаар хэрэглэгчдэд илүү ойлгомжтой, хүртээмжтэй, харилцан ярианд суурилсан үйлчилгээ үзүүлэх боломжтой.

Иймээс энэхүү дипломын ажил нь хүйсийн болон нийгмийн тэгш, хүртээмжтэй байдлыг хангахад чиглэсэн хиймэл оюун ухаант чатбот системийг судалж, түүний архитектур, мэдлэгийн сан, баримтад тулгуурласан хариулт үүсгэх аргачлалд суурилсан шийдлийг боловсруулахад чиглэж байна.



